\documentclass[a4paper,fleqn]{cas-sc}

% =========================================================================
% Academic Student Writer - Journal Paper Starter Template
% Layout Style: Elsevier CAS Single-Column (cas-sc) / Adaptable Standard
% Calibrated to 0% AI-Detection Benchmark (turnitin-verified authentic voice)
% =========================================================================

\usepackage{float}
\usepackage[utf8]{inputenc}
\usepackage[numbers]{natbib}
\usepackage{tabularx}
\usepackage[english]{babel}
\usepackage{amsmath}
\usepackage{amsfonts}
\usepackage{amssymb}
\usepackage{booktabs}
\usepackage{multirow}
\usepackage{hyperref}
\usepackage{subcaption}
\usepackage{algorithm}
\usepackage{algpseudocode}
\usepackage{color}
\usepackage{xcolor}
\usepackage{orcidlink}
\usepackage{lineno}
\usepackage{silence}
\WarningFilter{hyperref}{Ignoring empty anchor}

% Float placement parameters: prevent awkward page breaks with multiple floats
\setcounter{topnumber}{3}
\setcounter{bottomnumber}{2}
\setcounter{totalnumber}{5}
\renewcommand{\topfraction}{0.95}
\renewcommand{\bottomfraction}{0.95}
\renewcommand{\textfraction}{0.05}
\renewcommand{\floatpagefraction}{0.7}

% Tighter float and caption spacing standard in peer-reviewed journals
\AtBeginDocument{%
  \setlength{\textfloatsep}{10pt plus 2pt minus 2pt}%
  \setlength{\floatsep}{8pt plus 2pt minus 2pt}%
  \setlength{\intextsep}{8pt plus 2pt minus 2pt}%
  \setlength{\abovecaptionskip}{6pt}%
  \setlength{\belowcaptionskip}{2pt}%
}

\setlength{\emergencystretch}{2em}
\renewcommand{\arraystretch}{1.08}
\setlength{\tabcolsep}{5pt}

% Preprint footer definition
\ExplSyntaxOn
\cs_set:Npn \__first_footerline:
  {
    \group_begin:
    \small
    \sffamily
    \ifnum\theblind>0\relax
    \else
    \__short_authors: :~
    \fi
    \rmfamily \itshape Preprint Under Review
    \group_end:
  }
\bool_gset_true:N \g_stm_nologo_bool
\ExplSyntaxOff

% Graphicspath: Point to high-resolution figures folder (600+ DPI required)
\graphicspath{{./figures/}{figures/}}

\begin{document}
\let\WriteBookmarks\relax

\shorttitle{[Short Running Title]}
\shortauthors{[FirstAuthor et~al.]}

\title [mode = title]{[Full Manuscript Title: Specific, Empirical, and Grounded in Findings]}

% -------------------------------------------------------------------------
% Authors & Affiliations
% -------------------------------------------------------------------------
\author[1]{[Author One Name]}[orcid=[0000-0000-0000-0000]]
\cormark[1]
\ead{author1@institution.edu}
\credit{Conceptualization, Methodology, Software, Formal analysis, Investigation, Writing - original draft}

\author[1]{[Author Two Name]}[orcid=[0000-0000-0000-0000]]
\ead{author2@institution.edu}
\credit{Data curation, Validation, Writing - review \& editing}

\author[2]{[Author Three Name]}[orcid=[0000-0000-0000-0000]]
\ead{author3@institution.edu}
\credit{Supervision, Resources, Writing - review \& editing}

\affiliation[1]{organization={[Department of Computer Science and Engineering, University Name]},
  city={[City]},
  postcode={[Zip]},
  country={[Country]}}

\affiliation[2]{organization={[Department or Research Center, Institute Name]},
  city={[City]},
  postcode={[Zip]},
  country={[Country]}}

\cortext[1]{Corresponding author.}

% -------------------------------------------------------------------------
% Abstract (1 substantive paragraph, 150-250 words, concrete figures, no generic fluff)
% -------------------------------------------------------------------------
\begin{abstract}
[State the exact empirical problem in the first sentence with scope. Name the primary mechanism or model proposed. Specify the datasets and experimental baselines evaluated against. Report concrete performance numbers with exact deltas, error rates, or confidence intervals. Conclude with the practical operational implication or boundary condition of the findings, avoiding meta-claims such as ``paves the way'' or ``proves groundbreaking''.]
\end{abstract}

\begin{keywords}
[Keyword 1] \sep [Keyword 2] \sep [Keyword 3] \sep [Keyword 4] \sep [Keyword 5]
\end{keywords}

\maketitle

% =========================================================================
% SECTION ARCHITECTURE OVERVIEW (MAX 6 MAIN SECTIONS STANDARD)
%
% Choose the section skeleton matching your study type below:
%
% VARIANT A (Method-Focused Paper):
%   1. Introduction
%   2. Conventional Method
%   3. Proposed Method
%   4. Experiments
%   5. Results
%   6. Conclusion
%
% VARIANT B (Dataset-Focused Paper):
%   1. Introduction
%   2. Corpus Construction and Annotation
%   3. Baseline Benchmarks and Verification
%   4. Experiments and Validation Protocol
%   5. Results and Analysis
%   6. Conclusion
%
% VARIANT C (Mixed Method + Dataset Paper):
%   1. Introduction
%   2. Dataset and Task Formulation
%   3. Proposed Methodology
%   4. Experiments
%   5. Results
%   6. Conclusion
%
% RULE: Subsections must be limited to 4–6 core thematic groupings per section.
% Strictly avoid subsection sprawl (e.g., 4.1 to 4.8 or 5.1 to 5.10).
% Substantive paragraphs must be 150–250+ words; never fragment into 3–4 line micros.
% =========================================================================

\section{Introduction}\label{sec:intro}

[Paragraph 1: 180--240 words. Ground the real-world operational domain immediately with 3--4 relevant cited papers. Establish the concrete technical failure mode or friction point. Anchor specific numbers, hardware limits, or empirical discrepancies from prior literature.]

[Paragraph 2: 180--250 words. Analyze why existing literature (5--6 relevant prior studies) fails to resolve the friction point. Trace the exact structural trade-offs between speed, accuracy, calibration, or label efficiency. Explain the precise boundary where conventional mechanisms break down.]

[Paragraph 3: 180--240 words. Introduce the proposed approach by naming its core mathematical mechanism and structural guarantees. State clearly what is retained from conventional pipelines and what is modified or introduced. Enumerate the core empirical contributions plainly.]

% -------------------------------------------------------------------------
% SECTION 2: PRIOR FOUNDATIONS / CONVENTIONAL SETUP / DATASET FORMULATION
% -------------------------------------------------------------------------
\section{Conventional Method}\label{sec:conventional}
% For Dataset-Focused papers: \section{Corpus Construction and Annotation}
% For Mixed papers: \section{Dataset and Task Formulation}

\subsection{Preliminaries}
[Define the problem formally. Establish mathematical notation and state spaces. Define input features, target distributions, and loss objectives clearly.]

\subsection{Baseline Formulation}
[Detail the standard conventional pipeline currently in use. Explain its operational stages, feature extraction steps, and decision rules. Highlight the specific architectural assumptions that cause degradation under distribution or task shift.]

% -------------------------------------------------------------------------
% SECTION 3: PROPOSED METHODOLOGY
% -------------------------------------------------------------------------
\section{Proposed Method}\label{sec:proposed}

\subsection{Working Principle}
[Detail the conceptual design and workflow of the proposed system. Reference an end-to-end architecture diagram immediately.]

% -------------------------------------------------------------------------
% FIGURE GUIDELINES:
% 1. Render all figures at 600+ DPI (vector PDF preferred, or high-res PNG/TIFF).
% 2. All text within the infographic or chart must be large and legible (>= 8 pt).
% 3. Captions must be concise (1-2 sentences, <= 35 words); place all analysis in body prose.
% -------------------------------------------------------------------------
\begin{figure}[t]
  \centering
  % \includegraphics[width=0.95\textwidth]{pipeline.pdf}
  \caption{End-to-end pipeline and architectural workflow of the proposed framework.}
  \label{fig:pipeline}
\end{figure}

\subsection{Mathematical Formulation}
[Formulate the model components mathematically. Follow the strict rule: every display equation must be immediately followed by a lowercase, unindented ``where'' clause defining all newly introduced variables. Never write informal definitions like ``here a = ...''.]

The penalized split gain $G'$ balances empirical loss reduction against distribution divergence:
\begin{equation}
  G' = G - \lambda P(x), \label{eq:penalized_gain}
\end{equation}
where $G$ denotes the unpenalized gradient split gain, $P(x)$ is the non-negative penalty computed across evidence partitions, and $\lambda \in [0, 1]$ represents the regularization coefficient. When $P(x) = 0$, Equation~\eqref{eq:penalized_gain} reduces exactly to the conventional gradient formulation.

\subsection{Optimization Algorithm}
[Present the operational algorithm using the formal algorithm environment.]

\begin{algorithm}[H]
  \caption{Proposed model training and threshold calibration procedure}
  \label{alg:training}
  \begin{algorithmic}[1]
    \Require Training set $\mathcal{D}_{\text{train}}$, validation split $\mathcal{D}_{\text{val}}$, learning rate $\eta$, seeds $\mathcal{S}$
    \Ensure Calibrated model weights $\theta^*$, operational threshold $\tau^*$
    \State Extract feature representations for all input tuples in $\mathcal{D}_{\text{train}}$
    \State Apply monotone constraints to monotonic feature indices
    \For{each seed $s \in \mathcal{S}$}
      \State Initialize tree ensemble with seed $s$
      \State Optimize tree split objective defined in Equation~\eqref{eq:penalized_gain}
      \State Fit secondary probability calibrator on $\mathcal{D}_{\text{val}}$
      \State Compute false-positive budget threshold $\tau_s$
    \EndFor
    \State \Return best checkpoint $\theta^*$ and optimal operational threshold $\tau^*$
  \end{algorithmic}
\end{algorithm}

% -------------------------------------------------------------------------
% SECTION 4: EXPERIMENTS
% -------------------------------------------------------------------------
\section{Experiments}\label{sec:experiments}

\subsection{Experimental Setup}
[Describe the compute infrastructure (GPUs, CPU cores, memory), execution environment, library versions, fixed random seeds, and precision modes used for training and inference.]

\subsection{Datasets}
[Provide concrete dataset characteristics: sample counts, domain splits, label prevalence, and preprocessing steps. Summarize in a structured table.]

\begin{table}[b]
  \centering
  \caption{Summary of experimental benchmark datasets and split distributions.}
  \label{tab:datasets}
  \begin{tabular}{lrrrr}
    \toprule
    \textbf{Corpus} & \textbf{Total Rows} & \textbf{Train} & \textbf{Validation} & \textbf{Test} \\
    \midrule
    Benchmark-A     & 10{,}000            & 7{,}000        & 1{,}500             & 1{,}500       \\
    Benchmark-B     & 5{,}000             & 3{,}500        & 750                 & 750           \\
    \bottomrule
  \end{tabular}
\end{table}

\subsection{Baselines and Protocols}
[Identify all comparative models: classical statistical baselines, modern specialized models, and competitive LLM or foundation backbones. State exact hyperparameter settings and search protocols.]

\subsection{Evaluation Metrics}
[Provide complete, explicit mathematical definitions for all evaluation metrics used in the study. Never list metrics as undefined acronyms.]

To assess classification accuracy alongside probability calibration, models are evaluated using Macro-$F_1$, Area Under the Receiver Operating Characteristic curve (AUROC), and Expected Calibration Error (ECE). Macro-$F_1$ is defined as the harmonic mean of precision and recall:
\begin{equation}
  F_1 = \frac{2 \cdot \text{Precision} \cdot \text{Recall}}{\text{Precision} + \text{Recall}}, \quad
  \text{where } \text{Precision} = \frac{\text{TP}}{\text{TP} + \text{FP}}, \quad \text{Recall} = \frac{\text{TP}}{\text{TP} + \text{FN}},
\end{equation}
where $\text{TP}$, $\text{FP}$, and $\text{FN}$ denote true positives, false positives, and false negatives, respectively. For calibration, predictions are partitioned into $M$ equally spaced confidence bins $\{B_m\}_{m=1}^M$, and ECE measures the weighted absolute difference between accuracy and confidence:
\begin{equation}
  \text{ECE} = \sum_{m=1}^{M} \frac{|B_m|}{N} \left| \text{acc}(B_m) - \text{conf}(B_m) \right|,
\end{equation}
where $N$ is the total sample size, $\text{acc}(B_m)$ denotes empirical accuracy in bin $B_m$, and $\text{conf}(B_m)$ represents average predicted confidence in bin $B_m$.

% -------------------------------------------------------------------------
% SECTION 5: RESULTS AND DISCUSSION
% -------------------------------------------------------------------------
\section{Results}\label{sec:results}

\subsection{Model Comparison}
[Report primary test results across all benchmarks and comparative baselines. Report exact mean scores, standard deviations over seeds, and relative gains. Connect each result directly to the structural features of the models.]

\subsection{Component Ablation}
[Systematically isolate and remove individual components (e.g., penalty terms, specific feature blocks, constraint sets). Demonstrate the quantitative contribution of each module to final performance.]

\subsection{Calibration and Distribution Shift}
[Evaluate model performance under domain, task, or prompt shift. Include reliability diagrams, coverage statistics, or Brier scores to demonstrate calibration stability.]

\subsection{Statistical Analysis}
[Conduct rigorous hypothesis testing (e.g., McNemar's paired test, Wilcoxon signed-rank test, paired $t$-tests) with explicit $p$-values, effect sizes, and multiple-testing corrections such as Holm--Bonferroni. Never use ``significant'' colloquially without supporting statistical test figures.]

\subsection{Error Analysis and Limitations}
[Analyze representative failure cases where the model faltered. Dissect length bias, annotation noise, or boundary conditions. Discuss operational tradeoffs, computational latency, and resource costs.]

% -------------------------------------------------------------------------
% SECTION 6: CONCLUSION (STRICTLY 1 COHESIVE PARAGRAPH, 120-160 WORDS)
% -------------------------------------------------------------------------
\section{Conclusion}\label{sec:conclusion}
[Write exactly ONE cohesive paragraph of 120–160 words. Summarize the technical problem addressed, the specific architectural mechanism proposed, the verified empirical benchmark gains over baselines with exact figures, and the concrete operational boundary condition of the work. Do NOT start with ``In conclusion'' or ``Overall''. Do NOT write generic future work fluff (``paves the way'', ``exciting possibilities''). Ground the closure strictly in verified facts.]

% -------------------------------------------------------------------------
% Unnumbered Declarations & Acknowledgements
% -------------------------------------------------------------------------
\section*{Declarations}
\textbf{Funding:} [Funding details or ``No specific funding was received for this study.''] \\
\textbf{Conflict of interest:} The authors declare that they have no known competing financial interests or personal relationships that could have appeared to influence the work reported in this paper. \\
\textbf{Data availability:} [Code, model weights, and benchmark data will be made available at: URL / on request.]

\section*{Acknowledgements}
[Acknowledge institutional support, computing infrastructure, or advisory feedback.]

% -------------------------------------------------------------------------
% References
% -------------------------------------------------------------------------
\bibliographystyle{cas-model2-names}
\bibliography{references}

\end{document}
